\section{Task 7: Maximising the Quality}
\label{sec:task7}

\subsection{New Objective}
\label{sec:task7-objective}
To maximise quality, introduce a scalar variable $q$ representing the minimum achieved quality over the malt-source, hop, and yeast groups. The objective is
\begin{equation}
  \max q,
  \label{eq:max-quality-objective}
\end{equation}
or equivalently $\min -q$ when using \texttt{scipy.optimize.linprog}.

\subsection{Quality Constraints}
\label{sec:task7-constraints}
The volume, colour, and process constraints remain active. The three group-quality constraints are rewritten with $q$ as the lower bound:
\begin{align}
  \sum_{i \in B} Q^B_i(0.6b_i)
  +\sum_{i \in M} Q^M_i(0.8m_i)
  +\sum_{i \in E} Q^E_i e_i
  &\geq q\Mtot, \label{eq:q-source}\\
  \sum_{i \in H} Q^H_i h_i
  &\geq q\sum_{i \in H}h_i, \label{eq:q-hop}\\
  \sum_{i \in Y} Q^Y_i y_i
  &\geq q\sum_{i \in Y}y_i. \label{eq:q-yeast}
\end{align}
The bound $0\leq q\leq5$ is added. The binary process variables are still handled by the same brute-force enumeration as in Section~\ref{sec:relaxation-bruteforce}.

\subsection{Can the Minimum Quality Be 5?}
\label{sec:task7-quality5}
Yes. The optimisation reaches $q^\star=5.000$, so minimum quality 5 is feasible. Since 5 is also the upper bound on $q$, this is the best possible value. No ingredient group prevents quality 5 in the computed solution: the malt-source, hop, and yeast qualities all reach 5 while satisfying the amber colour constraint.

\begin{table}[!t]
\renewcommand{\arraystretch}{1.15}
\caption{Task 7 Quality Maximisation Result}
\label{tab:task37-quality}
\centering
\small
\begin{tabular}{@{}ll@{}}
\toprule
Item & Result \\
\midrule
Maximum minimum quality & 5.000 \\
Quality 5 feasible? & yes \\
Process choice & $(0,1)$ \\
Achieved EBC & 20.000 \\
Malt-source quality & 5.000 \\
Hop quality & 5.000 \\
Yeast quality & 5.000 \\
\bottomrule
\end{tabular}
\end{table}

\begin{table}[!t]
\renewcommand{\arraystretch}{1.15}
\caption{Task 7 Non-Zero Ingredients}
\label{tab:task37-ingredients}
\centering
\small
\begin{tabular}{@{}llr@{}}
\toprule
Group & Variable & Amount (kg) \\
\midrule
Malt & \texttt{malt11} & 870.772 \\
Malt extract & \texttt{maltextract03} & 66.541 \\
Hops & \texttt{hop10} & 6.842 \\
Yeast & \texttt{yeast01} & 3.750 \\
\bottomrule
\end{tabular}
\end{table}

% Values in this section come from task_37.ipynb.
